% ============================================================
% ANEXO A · CÓDIGO DE LOS MÓDULOS
% Fragmentos representativos de la implementación en notebooks.
% Decisiones que lo respaldan: D-008, D-014, D-017, D-018, D-020,
% D-021, D-022, D-023, D-024, D-025, D-026, D-027, D-031, D-032,
% D-033, D-036, D-039, D-040.
% ============================================================

\section{Código de los módulos}\label{anx:codigo}

En este anexo se presentan los fragmentos de código más representativos del pipeline
experimental desarrollado en el proyecto. Todos los listados incluidos han sido extraídos
de manera literal directamente desde las celdas ejecutadas de los cuadernos de trabajo
(\textit{Jupyter notebooks}) mediante la herramienta automatizada del repositorio.
Las rutinas auxiliares de utilidad general localizadas en el módulo \texttt{Code/tools/}
son omitidas en los listados para centrar la atención en la lógica nuclear de cada etapa.

% ------------------------------------------------------------
\subsection{Módulo 2: Preprocesamiento y alineación de máscaras}\label{sec:anx_m2}

% fuente: D-020
La función \texttt{align\_mask} implementa la regla de corrección espacial para las 80
máscaras del CBIS-DDSM cuyas dimensiones difieren de la mamografía completa. Cuando el
desajuste se debe a un factor de submuestreo de 0.870 se aplica interpolación al vecino
más cercano, mientras que en discrepancias de resolución nativa se preserva el anclaje
en la esquina superior izquierda sin remuestreo continuo.

\lstinputlisting[style=tt, caption={[Función de alineación espacial de máscaras]Función de alineación espacial de máscaras (\texttt{Code/2\_Preprocessing.ipynb}, celda de código 5).}, label={lst:m2_alineacion}]{codigo/m2_alineacion.py}

% fuente: D-021
La rutina \texttt{prepare\_mask} realiza el análisis morfológico de componentes conexas
mediante la función \texttt{connectedComponentsWithStats} de OpenCV. Identifica la componente de mayor área
como la lesión patológica principal y descarta motas satélite aisladas, registrando la
distancia euclidiana máxima y el número de píxeles eliminados para control de calidad.

\lstinputlisting[style=tt, caption={[Depuración y selección de componente conexa mayor]Depuración y selección de componente conexa mayor (\texttt{Code/2\_Preprocessing.ipynb}, celda de código 11).}, label={lst:m2_componente}]{codigo/m2_componente.py}

% fuente: D-022
La función \texttt{export\_mammogram} formaliza el contrato de salida del Módulo 2. Calcula
la caja delimitadora ajustada a la lesión depurada y le agrega un margen de seguridad de 20
píxeles por lado, extrayendo los parches sobre las intensidades crudas de 16 bits y
almacenándolos en formato comprimido NRRD con espaciado espacial unitario.

\lstinputlisting[style=tt, caption={[Recorte de región de interés con margen y exportación]Recorte de región de interés con margen y exportación (\texttt{Code/2\_Preprocessing.ipynb}, celda de código 13).}, label={lst:m2_recorte}]{codigo/m2_recorte.py}

% ------------------------------------------------------------
\subsection{Módulo 3: Extracción de características radiómicas}\label{sec:anx_m3}

% fuente: D-008, D-023
La inicialización del extractor radiómico fija el parámetro \texttt{binCount=32} sobre las
intensidades originales, asegurando invarianza a traslaciones lineales sin forzar un ancho
fijo de contenedor. Se inhabilitan los filtros derivados complejos y se activan únicamente
las cuatro familias morfológicas y texturales bidimensionales primarias (\textit{firstorder},
\textit{shape2D}, \textit{glcm} y \textit{glrlm}), totalizando 67 características.

\lstinputlisting[style=tt, caption={[Configuración del extractor PyRadiomics y activación de familias]Configuración del extractor PyRadiomics y activación de familias (\texttt{Code/3\_Feature\_Extraction.ipynb}, celda de código 3).}, label={lst:m3_extractor}]{codigo/m3_configuracion.py}

% ------------------------------------------------------------
\subsection{Módulo 4: Selección y escalado angular}\label{sec:anx_m4}

% fuente: D-024
La función \texttt{select} aplica un procedimiento voraz guiado por el estadístico $F$ del
análisis de varianza univariado, calculado estrictamente sobre los datos de entrenamiento.
Cada variable candidata se incorpora solo si el valor absoluto de su correlación de Pearson
frente a las variables ya seleccionadas no excede el límite de redundancia $|r| \le 0.95$.

\lstinputlisting[style=tt, caption={[Selección voraz por $F$-test con restricción de redundancia lineal]Selección voraz por $F$-test con restricción de redundancia lineal (\texttt{Code/4\_Selection\_and\_Scaling.ipynb}, celda de código 5).}, label={lst:m4_seleccion}]{codigo/m4_seleccion.py}

% fuente: D-025
El escalado angular utiliza \texttt{QuantileTransformer} configurado para proyectar a una
distribución uniforme y multiplicado por $\pi$. Esto asegura que cada característica ocupe
de forma homogénea el intervalo angular $[0, \pi]$, mitigando la influencia de asimetrías
y valores atípicos sin alterar el orden relativo de las observaciones de entrenamiento.

\lstinputlisting[style=tt, caption={[Transformación angular por cuantiles acotada a {$[0, \pi]$}]Transformación angular por cuantiles acotada a {$[0, \pi]$} (\texttt{Code/4\_Selection\_and\_Scaling.ipynb}, celda de código 8).}, label={lst:m4_escalado}]{codigo/m4_escalado.py}

% fuente: D-018, D-026
La conformación de los pliegues de validación cruzada utiliza \texttt{StratifiedGroupKFold}
para agrupar las lesiones por paciente. La rutina evalúa múltiples semillas deterministas y
selecciona aquella que minimiza la desviación máxima en la proporción de malignidad entre los
cinco pliegues, garantizando particiones homogéneas e independientes a nivel de sujeto.

\lstinputlisting[style=tt, caption={[Generación y balanceo de folds agrupados por paciente]Generación y balanceo de folds agrupados por paciente (\texttt{Code/4\_Selection\_and\_Scaling.ipynb}, celda de código 9).}, label={lst:m4_folds}]{codigo/m4_folds.py}

% ------------------------------------------------------------
\subsection{Módulo 5: Transformaciones cuánticas y kernels}\label{sec:anx_m5}

% fuente: D-014, D-031
El cálculo del embedding cuántico combina los circuitos de codificación C4 y C5 con un
\textit{ansatz} variacional \texttt{RealAmplitudes} parametrizado mediante un vector $\theta$
fijo generado con semilla 42. El valor esperado local $\langle Z_i \rangle$ se obtiene de forma
exacta multiplicando la distribución de probabilidades del vector de estado por la matriz de signos.

\lstinputlisting[style=tt, caption={[Construcción de mapas cuánticos, ansatz con $\theta$ fijo y cálculo de $\langle Z_i \rangle$]Construcción de mapas cuánticos, ansatz con $\theta$ fijo y cálculo de $\langle Z_i \rangle$ (\texttt{Code/5a\_Embeddings.ipynb}, celda de código 2).}, label={lst:m5_embeddings}]{codigo/m5_embeddings.py}

% fuente: D-017, D-027
La computación de matrices de kernel emplea \texttt{FidelityStatevectorKernel} sobre los
vectores cuánticos precomputados, optimizando el cálculo de productos internos exactos en
la submuestra representativa. En paralelo, se construye el kernel clásico RBF empleando el
ancho de banda óptimo determinado por la heurística de la mediana.

\lstinputlisting[style=tt, caption={[Evaluación de matrices de kernel clásico y cuántico]Evaluación de matrices de kernel clásico y cuántico (\texttt{Code/5b\_Quantum\_Kernels.ipynb}, celda de código 4).}, label={lst:m5_kernels}]{codigo/m5_kernels.py}

% ------------------------------------------------------------
\subsection{Módulo 6: Separabilidad geométrica}\label{sec:anx_m6}

% fuente: D-033
El alineamiento de matrices de kernel evalúa la correspondencia con la estructura de clases.
La rutina \texttt{alignment\_with\_null} centra la matriz de kernel y calcula tanto la
métrica empírica observada como su exceso respecto a la media de una distribución nula
obtenida mediante permutaciones exhaustivas de las etiquetas binarias.

\lstinputlisting[style=tt, caption={[Centrado matricial y alineamiento kernel-objetivo con prueba de permutación]Centrado matricial y alineamiento kernel-objetivo con prueba de permutación (\texttt{Code/6\_Separability.ipynb}, celda de código 3).}, label={lst:m6_kta}]{codigo/m6_kta_permutaciones.py}

% fuente: D-017, D-032
La evaluación de ventaja cuántica potencial implementa la diferencia geométrica regularizada
$g(K_C \parallel K_Q)$ descrita por Huang et al. El algoritmo computa la raíz simétrica
de la matriz cuántica y diagonaliza el kernel clásico, evaluando las cotas teóricas de
generalización y entrenamiento sobre una grilla de parámetros de regularización $\lambda$.

\lstinputlisting[style=tt, caption={[Cálculo de la diferencia geométrica regularizada de Huang et al]Cálculo de la diferencia geométrica regularizada de Huang et al. (\texttt{Code/6\_Separability.ipynb}, celda de código 3).}, label={lst:m6_geometric_difference}]{codigo/m6_geometric_difference.py}

% ------------------------------------------------------------
\subsection{Módulo 7: Clasificación supervisada}\label{sec:anx_m7}

% fuente: D-036
El clasificador neuronal común consiste en un perceptrón multicapa implementado en PyTorch
con estructura $k \to 32 \to 16 \to 1$, capas lineales continuas con activaciones ReLU y
salida en escala logit, lo que permite integrar numéricamente la función sigmoide dentro de la
función de costo de entropía cruzada binaria ponderada.

\lstinputlisting[style=tt, caption={[Definición del perceptrón multicapa para clasificación supervisada]Definición del perceptrón multicapa para clasificación supervisada (\texttt{Code/7\_Classification.ipynb}, celda de código 3).}, label={lst:m7_mlp}]{codigo/m7_mlp.py}

% fuente: D-036
El ciclo de optimización \texttt{fit\_predict} utiliza el optimizador Adam y aplica parada
temprana monitoreando la pérdida ponderada sobre una partición interna independiente de
pacientes. Si tras 40 épocas no se detecta mejoría, se detiene el entrenamiento y se
restauran los mejores pesos obtenidos para evaluar el conjunto de prueba.

\lstinputlisting[style=tt, caption={[Bucle de entrenamiento con parada temprana y restauración de pesos]Bucle de entrenamiento con parada temprana y restauración de pesos (\texttt{Code/7\_Classification.ipynb}, celda de código 3).}, label={lst:m7_parada_temprana}]{codigo/m7_parada_temprana.py}

% ------------------------------------------------------------
\subsection{Módulo 8: Evaluación comparativa y muestreo}\label{sec:anx_m8}

% fuente: D-039
El análisis comparativo evalúa la correlación de rangos de Spearman entre las métricas de
separabilidad geométrica y el desempeño clasificador. La rutina asigna rangos promedio a
condiciones empatadas y deriva la distribución nula exacta mediante las 120 permutaciones
posibles para determinar la significación estadística bilateral.

\lstinputlisting[style=tt, caption={[Cálculo del coeficiente de Spearman con empates y distribución nula exacta]Cálculo del coeficiente de Spearman con empates y distribución nula exacta (\texttt{Code/8\_Comparison.ipynb}, celda de código 6).}, label={lst:m8_spearman}]{codigo/m8_spearman.py}

% fuente: D-040
La cuantificación del coste por mediciones finitas simula el comportamiento de un dispositivo
físico mediante muestreo multinomial de las probabilidades exactas. Este procedimiento emula
un ejecutor ideal de disparos finitos (\textit{shots}) y permite contrastar el error empírico
frente a la cota teórica de desviación estándar de la medición cuántica.

\lstinputlisting[style=tt, caption={[Muestreo multinomial de mediciones finitas y ley de dispersión]Muestreo multinomial de mediciones finitas y ley de dispersión (\texttt{Code/8\_Comparison.ipynb}, celda de código 17).}, label={lst:m8_muestreo}]{codigo/m8_muestreo.py}
